\documentclass[12pt]{article}
\usepackage{hyperref}

\usepackage{graphicx}

\usepackage{mathtools}
\DeclarePairedDelimiter\ceil{\lceil}{\rceil}
\DeclarePairedDelimiter\floor{\lfloor}{\rfloor}


\usepackage{amsmath}
\usepackage{amsfonts}
\usepackage{amssymb}

\usepackage{bm}

\newcommand{\Mod}[1]{\ (\mathrm{mod}\ #1)}


\usepackage{listings}
\lstset{language=C++,
                basicstyle=\ttfamily,
                keywordstyle=\color{blue}\ttfamily,
                stringstyle=\color{red}\ttfamily,
                commentstyle=\color{green}\ttfamily,
                morecomment=[l][\color{magenta}]{\#}
}

\usepackage{breqn}
\usepackage[]{algorithm2e}
\usepackage{physics}

\newcommand\fnurl[2]{%
  \href{#2}{#1}\footnote{\url{#2}}%
}

\newcommand{\Four}{\mathcal{F}}
\renewcommand{\(}{\left(}
\renewcommand{\)}{\right)}
%% \renewcommand{\{}{\left{}
%% \renewcommand{\}}{\right}}


\title{Convolution Notes for rocFFT}
\author{The rocFFT Team}



\def\no#1{\nobreak{#1}} % stop equation breaking in dmath

\begin{document}
\maketitle

\section{Introduction}


The convolution of two inputs $\{f\}$ and $\{g\}$ of equal length $N$ is
\begin{dmath}
  \(f\star{}g\)_k = \sum_{p=0}^{k}f_{k-p}g_p
\end{dmath}
for $n=0,\dots,N-1$.  The computational complexity is $O\(N^2\)$.

This can be extended to $d$ dimensions, where the multi-index'd convolution is
\begin{dmath}
  \(f\star{}g\)_{\pmb{n}} = \sum_{m_0=0}^{k_0}\dots{}\sum_{m_{d-1}=0}^{k_{d-1}} f_{\pmb{n}-\pmb{m}}g_{\pmb{m}}
\end{dmath}
TODO: what's the complexity here?

\section{Convolution via FFT}

We need just one index with $-k +p +q$ modulo the length of the FFT.

For $k=0$ we should only have one contribution to this sum.
\begin{dmath}
  k  = -k + p + q = 0 + p + q
\end{dmath}
where the equality is modulo the FFT length.  Clearly, $p=q=0$ is a
solution.  The first non-zero solution, if it exists, would be
achieved at $p=q=N-1$, making the sum equal to $2N-2$.  In order to
suppress this second solution, the FFT must be of length at least
$2N-1$.

    

\end{document}
